% TOP_PROBLEM: {"id": 20001072, "title": "Exact feedback sets in cyclic tournament substitutions", "category_id": 2, "category": {"id": 2, "name": "combinatorics", "display_name": "Combinatorics", "description": "Counting problems, graph theory, discrete structures.", "slug": "combinatorics", "order_index": 2, "created_at": "2026-07-31T15:26:25.671Z"}, "status": "partially_solved", "problem_number": "AIM-COMBINATORICS-0197", "set_id": 14}
% FIRST_POSTED: 2026-10-01
% ENTRY_KIND: statement_only
% UnsolvedMath ID 20001072; source catalogue code AIM-COMBINATORICS-0197
% !TeX program = lualatex
% Definition and sources only; this file is not an LLM solution attempt.
\documentclass[11pt,a4paper]{article}
\usepackage[margin=25mm]{geometry}
\usepackage{fontspec}
\setmainfont{lmroman10-regular.otf}[BoldFont=lmroman10-bold.otf,
 ItalicFont=lmroman10-italic.otf,BoldItalicFont=lmroman10-bolditalic.otf]
\usepackage{amsmath,amssymb,mathtools}
\usepackage{unicode-math}
\setmathfont{latinmodern-math.otf}
\AtBeginDocument{\let\setminus\smallsetminus}
\usepackage{xurl,hyperref}
\hypersetup{colorlinks=true,urlcolor=blue}
\setcounter{secnumdepth}{0}
\setlength{\parindent}{0pt}
\setlength{\parskip}{5pt}
\setlength{\emergencystretch}{3em}
\begin{document}
\section*{Exact feedback sets in cyclic tournament substitutions}\label{20001072}
{\small UnsolvedMath ID 20001072 · AIM-COMBINATORICS-0197 · Combinatorics · AIM Workshop Problem Lists}
\subsection{Definitions and mathematical statement}
Let $D=(V,E)$ be a finite simple digraph with no directed cycle of
length at most three. Thus loops, opposite arcs, and cyclically
oriented triangles are all absent. An unordered pair of distinct
vertices $\{u,v\}$ is a nonedge if neither $u\to v$ nor $v\to u$
belongs to $E$. Let
\[
             k=\left|\left\{\{u,v\}\in\binom V2:
                          u\to v,v\to u\notin E\right\}\right|.
\]
A feedback arc set is a subset $F\subseteq E$ such that deleting
its arcs leaves no directed cycle of any length. Write
$\tau(D)=\min\{|F|:F\text{ is a feedback arc set}\}$.

The Chudnovsky--Seymour--Sullivan conjecture asks whether
\[
                              \tau(D)\le\lfloor k/2\rfloor
\]
for every such digraph. Missing pairs are counted once, while deleted
arcs are individual directed edges. Vertices are retained throughout.
Equivalently, there should be a linear ordering of $V$ with at most
$\lfloor k/2\rfloor$ arcs pointing backwards: deleting these arcs
leaves an acyclic graph, and every finite acyclic digraph admits an
ordering in which all arcs point forwards.
The hypothesis only forbids short directed cycles, while the
conclusion removes all directed cycles. A transitive triangle is
permitted, and the original graph need not have any particular
ordering, connectedness, or degree conditions.
\subsection{Short English statement}
Can every digraph without a directed triangle or digon be made acyclic by deleting at most half as many arcs as it has missing unordered pairs?
\subsection{Sources}
\begin{itemize}
\item[S1] ULAM.AI and the UnsolvedMath Contributors, supplied problems.json, record 20001072 (AIM-COMBINATORICS-0197), AIM Workshop Problem Lists. Catalogue identity and source transcription; CC BY 4.0.
\url{https://huggingface.co/datasets/ulamai/UnsolvedMath}.
\item[S2] American Institute of Mathematics, The Caccetta-Haggkvist conjecture, item 6. Original workshop question underlying AIM-COMBINATORICS-0197.
\url{https://aimath.org/WWN/caccetta/caccetta.pdf}.
\end{itemize}
\end{document}
